\documentclass[10pt]{article}
\usepackage[T1]{fontenc}
\usepackage[utf8]{inputenc}
\usepackage[french]{babel}
\usepackage{amsmath,amssymb,mathrsfs,mathpazo}
\usepackage{geometry}
\geometry{a4paper,margin=2cm}
\usepackage{graphicx,array,tabularx,tabularray,colortbl,multirow,multicol,booktabs,enumitem,pifont,ulem,fancybox,diagbox,cancel}
\usepackage{pstricks,pstricks-add,pst-plot,pst-tree,pst-node,pst-text,pst-func,pst-math,pst-bspline,pst-3dplot,pst-eucl,pst-solides3d}
\let\pscalculate\undefined
\usepackage{pst-bezier}
\usepackage{tikz,pgfplots,tkz-tab}
\usetikzlibrary{arrows,arrows.meta,patterns,calc,positioning,trees,decorations.markings,shapes.geometric,shapes.misc,shapes.arrows,intersections,decorations.pathmorphing,decorations.pathreplacing,backgrounds,shadows,babel}
\pgfplotsset{compat=1.18}
\usepgfplotslibrary{fillbetween}
\usepackage[np]{numprint}
\usepackage{listings,makecell,pas-tableur,twoopt}
\let\barre\undefined
\usepackage[pstricks]{tablvar}
\definecolor{darkgreen}{RGB}{0,100,0}
\providecommand{\red}{\color{red}}
\providecommand{\blue}{\color{blue}}
\providecommand{\green}{\color{green}}
\providecommand{\black}{\color{black}}
\providecommand{\white}{\color{white}}
\providecommand{\gray}{\color{gray}}
\providecommand{\cyan}{\color{cyan}}
\providecommand{\magenta}{\color{magenta}}
\providecommand{\yellow}{\color{yellow}}
\providecommand{\violet}{\color{violet}}
\providecommand{\orange}{\color{orange}}
\providecommand{\euro}{\texteuro}
\providecommand{\EUR}{\texteuro}
\providecommand{\ieme}{\textsuperscript{e}}
\providecommand{\degres}{\ensuremath{^\circ}}
\providecommand{\degree}{\ensuremath{^\circ}}
\providecommand{\vv}[1]{\overrightarrow{#1}}
\providecommand{\Coord}[1]{#1}
\providecommand{\decosix}{}
\providecommand{\starredbullet}{\ensuremath{\bullet}}
\pagestyle{empty}
\setlength{\parindent}{0pt}
\usepackage{adjustbox}\DeclareUnicodeCharacter{2236}{:}\DeclareUnicodeCharacter{2190}{\ensuremath{\leftarrow}}
\providecommand{\R}{\mathbb{R}}
\providecommand{\N}{\mathbb{N}}
\providecommand{\Z}{\mathbb{Z}}
\providecommand{\Q}{\mathbb{Q}}
\providecommand{\D}{\mathbb{D}}
\providecommand{\C}{\mathbb{C}}
\providecommand{\e}{{\mathrm{e}}}
\providecommand{\vect}[1]{\overrightarrow{#1}}
\providecommand{\vectt}[1]{\overrightarrow{\mathrm{#1}}}
\providecommand{\vvt}[1]{\overrightarrow{\mathrm{#1}}}
\providecommand{\np}[2][]{{#2~#1}}
\providecommand{\vv}[1]{{\overrightarrow{#1}}}
\renewcommand{\d}{{\mathrm{d}}}
\providecommand{\degres}{{^\circ}}
\providecommand{\degre}{{^\circ}}
\providecommand{\blue}{\color{blue}}
\providecommand{\red}{\color{red}}
\providecommand{\enskip}{\quad}
\providecommand{\llbracket}{\lbrack\!\lbrack}
\providecommand{\rrbracket}{\rbrack\!\rbrack}
\providecommand{\euro}{{\text{€}}}
\providecommand{\no}{{\text{n°}}}
\providecommand{\interval}[2]{{[#1;#2]}}
\providecommand{\intervFO}[2]{{[#1;#2[}}
\providecommand{\dotfill}{\ldots}

\providecommand{\ds}{\displaystyle}
\providecommand{\textding}[1]{#1}
\providecommand{\Oij}{\(\mathrm{O};\vec{\imath},\vec{\jmath}\)}
\providecommand{\Oijk}{\(\mathrm{O};\vec{\imath},\vec{\jmath},\vec{k}\)}
\providecommand{\Ouv}{\(\mathrm{O};\vec{u},\vec{v}\)}
\providecommand{\intervFF}[2]{\ensuremath{[#1;#2]}}
\renewcommand{\euro}{\ifmmode\text{\texteuro}\else\texteuro\fi}
\DeclareUnicodeCharacter{00B0}{\ifmmode^{\circ}\else\textdegree\fi}
\DeclareUnicodeCharacter{20AC}{\ifmmode\text{\texteuro}\else\texteuro\fi}
% Énoncés : Éducation nationale. Transcription et corrigés : APMEP, Denis Vergès et contributeurs. Découpage et classement : CoopMaths, Sébastien Lozano, Éric Elter et contributeurs. Adaptation de la présentation : Maths Explorer.
% CC BY-SA : https://coopmaths.fr/www/about/
% https://coopmaths.fr/alea/static/eam/2026/tex/eam_specialite_2026_06_asie_automatismes_cor.tex
\begin{document}
\noindent\textit{APMEP / CoopMaths. CC BY-SA. Présentation adaptée par Maths Explorer.}\par\medskip
\subsection*{Question 1}
	La représentation graphique de $f$ a trois points d'intersection avec la droite d'équation {$y=-1$}, l'équation {$f(x)=-1$} a donc trois solutions.

	\textbf{réponse D}.

\subsection*{Question 2}

	Multiplier par {$\np{0.6}=\dfrac{60}{100}=\dfrac{100-40}{100}$} revient bien à diminuer la quantité de départ de 40\,\%.

	\textbf{réponse C}.

\subsection*{Question 3}

	La droite passe par les points {$(0;2)$} et {$(6;0)$}, donc son ordonnée à l'origine est $2$, et sa pente {$\dfrac{0-2}{6-0}=-\dfrac13$}.

	\textbf{réponse C}.

\subsection*{Question 4}

	Il y a 130 enfants, dont 80 pratiquent le volley-ball, donc {$P_E(V)=\dfrac{80}{130}=\dfrac{8}{13}$}.

	\textbf{réponse A}.

\subsection*{Question 5}

	$A$ et $\overline{A}$ partitionnent l'univers. D'après la loi des probabilités totales, on a :

	{$P(B)=P(A\cap B)+P(\overline{A}\cap B)=\np{0.4}\times\np{0.3}+\np{0.6}\times\np{0.1}=\np{0.12}+\np{0.06}=\np{0.18}$}.

	\textbf{réponse B}.

\subsection*{Question 6}
	Comme on sait que $f$ est une fonction affine, elle est représentée par une droite, et comme $g$ est une fonction polynôme du second degré, elle est représentée par une parabole, ce qui permet d'identifier les fonctions et leurs courbes représentatives.

	La droite représentant $f$ est au-dessous de la parabole représentant $g$ pour des valeurs de $x$ inférieures à $-2$ ou bien supérieures à $2$, donc la solution de l'inéquation {$f(x)<g(x)$} est {$S=\,]-\infty ; -2[\,\cup\,]2 ; +\infty[$}.

	\textbf{réponse B}.

\subsection*{Question 7}
	$x$ ne peut pas être égal à $3$ car la fonction inverse ne s'annule jamais.

	{$\begin{aligned}[t]
			x=3+\dfrac5 y
			&\iff x-3=\dfrac 5 y\\
			&\iff \dfrac{1}{x-3}=\dfrac y 5\\
			&\iff y=\dfrac{5}{x-3}
		\end{aligned}$}.

	\textbf{réponse D}.

\subsection*{Question 8}
	Additionnons les effectifs : \quad $3 + 1 + 2 + 2 + 2 = 10$.

	Il y a 10 valeurs, donc la médiane est entre les 5\up{e} et 6\up{e} valeurs, qui sont toutes deux 20 : en clair, $\textit{Me} =20 $.

	Ici, on n'a même pas besoin de calculer exactement la moyenne (mais on pourrait le faire, la division par 10 est simple).

	Toutes les valeurs sont positives et il y a deux valeurs qui valent 100, donc la somme de toutes les valeurs est strictement supérieure à 200, et la moyenne est donc strictement supérieure à 20.

	En conclusion, \quad {$\textit{Me} < \overline{x}$}.

	\textbf{réponse C}.

\end{document}
